% HD-I12. GPT-6 Astra Pro, inherited effort, /root/imprimitive.
% Written and author-read 21 November 2026.
\subsection{HD-I12: exact ranks on whole reflection branch strata}
\label{sec:HD-I12}

Fix $d\in\{3,4,6\}$, $e\mid d$, $n\geq5$, and
$m\geq\lceil(n^2-2n+3)/2\rceil$.
Consider a Nielsen component for the standard representation of
$G(d,e,m)$ all of whose branch generators are reflections.  A reflection
has exactly one of the following two forms:
\begin{enumerate}
\item a colored transposition $s_{ij}(a)$, with its two-cycle product equal
to $1$ and all other coordinates fixed;
\item a diagonal rotation $\operatorname{diag}(\zeta_d^s,1,\ldots,1)$,
up to permutation, with $0<s<d$ and $e\mid s$.
\end{enumerate}
Let $T$ be the number of generators of the first form, and $r_s$ the
number of the second form with color $s$.  Set $r=\sum_s r_s$.
The quotient tower of HD-I01--HD-I03 is
\[
 X=Z/\ker\chi\ \xrightarrow{\;\pi\;}\ Y=Z/H
 \ \xrightarrow{\;h\;}\ \mathbf P^1.
\]
Here $h$ is simply branched of degree $m$, its branch values are distinct,
and the $r$ branch points of $\pi$ have distinct images disjoint from
these values.  In particular
\begin{equation}
 g=1-m+\frac T2,\qquad
 r=2(n-g+1),\qquad \sum_s s r_s=d(n-g+1).
 \tag{I12.1}
\end{equation}
The actual polarized isogeny of HD-I02 identifies the multiplicity factor
with the primitive cyclic factor up to the fixed multiplier $|\ker\chi|$.
Consequently the following calculation is a rank calculation for the
original $G(d,e,m)$ period map.

\paragraph{Statement.}
For every such Nielsen component with $0\leq g\leq n-2$, the period map
has constant differential rank, on its smooth finite-level cover,
\begin{equation}
 \operatorname{rank}(d\mathrm{Per})=2n+g-1.
 \tag{I12.2}
\end{equation}
This determines each of these individual components; it is not a maximum
over components.  The full number of branch values is
$T+r=2m+2n$, so its Hurwitz dimension is $2m+2n-3$, and its period-fibre
dimension is exactly $2m-g-2$.

\paragraph{1. The cyclic multiplication map is surjective at every point.}
Write $R$ for the reduced branch divisor and choose the root
$L^d=\mathcal O_Y(\sum_s sR_s)$, as in HD-I05.  Define
\[
 A=\omega_Y L,\qquad B=\omega_Y L^{-1}(R).
\]
Equations (I12.1) give
\[
 \deg A=\deg B=n+g-1,\quad h^0(A)=h^0(B)=n,
 \quad AB=\omega_Y^2(R).
\]
For $g\leq n-2$, these two degrees are at least $2g+1$.  In particular
$A$ is globally generated: for any $p\in Y$, Serre duality gives
$H^1(A(-p))=H^0(\omega_Y A^{-1}(p))^*=0$, since the latter line bundle
has negative degree.  Riemann--Roch gives the two stated dimensions.
For the degree-zero line bundle $M=BA^{-1}$, a nonzero section either
trivializes $M$ or is impossible.  Thus $h^0(M)\leq1$ and
\[
 h^1(M)=g-1+h^0(M)\leq g\leq n-2.
\]
Green's $H^0$ lemma, with $W=H^0(A)$, now gives a surjection
\begin{equation}
 H^0(A)\otimes H^0(B)\longrightarrow H^0(AB).
 \tag{I12.3}
\end{equation}
The lemma says precisely that a base-point-free subspace
$W\subset H^0(A)$ has this surjectivity whenever
$h^1(BA^{-1})\leq\dim W-2$; there is no additional general-curve or
general-line-bundle hypothesis.  The original locator is M.~Green,
\emph{Koszul cohomology and the geometry of projective varieties},
J.~Differential Geom.~19 (1984), 125--171, Theorem~(4.e.1).
The exact statement was read in the published restatement by
B.~P.~Purnaprajna, \emph{Some Results on Surfaces of General Type},
Canad.~J.~Math.~57 (2005), 724--749, Lemma~2.7, printed p.~729,
\href{https://doi.org/10.4153/CJM-2005-029-3}{doi:10.4153/CJM-2005-029-3}.
The original Green PDF was not accessible in this check; this source
statement is therefore recorded as a checked published restatement,
not as a fresh reading of Green's original proof.

By HD-I05, (I12.3) is exactly the codifferential of the primitive cyclic
period map, up to the nonzero trace scalar $d$.  Since
$\deg(AB)=2n+2g-2>2g-2$, its target has dimension
\[
 h^0(\omega_Y^2(R))=2n+g-1.
\]
Thus that period differential is injective on the full pointed cyclic
cover deformation space, of this dimension.

\paragraph{2. Every infinitesimal cyclic deformation lifts inside the same
Nielsen component.}
Let $Q=h^*\mathcal O_{\mathbf P^1}(1)$, so $\deg Q=m$.
The obstruction to lifting the map $h$ over a deformation of $Y$,
with target $\mathbf P^1$ fixed, belongs to
\[
 H^1(Y,h^*T_{\mathbf P^1})=H^1(Y,Q^2)=0.
 \tag{I12.4}
\]
Indeed $\deg Q^2=2m>2g-2$.  The numerical range suffices even for
$g\leq n+1$: as proved in HD-I10, $m\geq n+4$.
Here is the obstruction map explicitly.  Trivialize a first-order
deformation of $Y$ on an affine cover.  Smoothness of $\mathbf P^1$
lifts each local map.  Differences of the local lifts on overlaps are
derivations with values in $h^*T_{\mathbf P^1}$; they form a Cech
$1$-cocycle.  Vanishing in (I12.4) lets one subtract local derivations
to make the lifts agree.  This works as well for a square-zero
extension with any finite-dimensional ideal, by tensoring the same
cohomology group with that ideal.

The marked branch points of $\pi$ are already part of the chosen cyclic
deformation.  They impose no fixed-image conditions on $h$.
Simple ramification of $h$, pairwise distinct critical values, and
separation of the images of the marked points are open conditions,
so they remain valid locally.  Roots $L^d$ of the moving weighted
divisor deform with the cyclic cover: multiplication by $d$ on the
relative Picard is etale, since the characteristic is zero.
The composite degree-$md$ cover and its Galois closure therefore deform
with unchanged topological monodromy, after a finite marking.  In
particular the local lift lies in the original Nielsen component,
not just in an unspecified component with the same signature.
Equivalently, forgetting $h$ from this tower has surjective differential
onto $H^1(T_Y(-R))$ at every point.

Combining this surjection, the injective cyclic period differential,
and the invertible differential of the fixed isogeny correspondence
in HD-I02 proves (I12.2).

\paragraph{3. Nonempty strata and explicit smaller-rank components.}
This calculation applies to every nonempty reflection stratum with
the data (I12.1).  The following construction supplies such strata,
including the concrete genus-zero examples.
Fix any balanced nonnegative counts $r_s$ with $r>0$, and put
$e_0=\gcd(d,\{s:r_s>0\})$.  For $g=0$ take $e_0=1$; for $g\geq1$
allow any $e_0\mid d$ with $e_0<d$.  On a smooth genus-$g$ curve choose
distinct marked points with these colors.  The cyclic character of the
punctured surface sends each small branch loop to its specified color.
The surface relation is satisfied because $\sum_s s r_s\equiv0\pmod d$.
For $g=0$ the images generate $\mathbf Z/d$ by $e_0=1$; for $g\geq1$
send one handle generator to $1$ and its paired generator to $0$.
This defines a connected cyclic cover with exactly the prescribed
local data.  Riemann existence makes it algebraic.

Choose a simply branched degree-$m$ pencil $h$ whose critical values and
marked images are all distinct.  Its existence follows from the
very-ample-pencil construction in HD-I08, since $m\geq g+3$ in the
present range.  The graph and Bezout proof of HD-I08 extends to the
several diagonal colors without any change of map: the diagonal
generators are now all $sE_i$, and their combinations give exactly
$e_0E_i$.  After the spanning-tree gauge, the extra colored edges give
$a(E_i-E_j)$.  If
$q=\gcd(d,e_0,\{a\})$, the generated diagonal subgroup is
\[
 \{v\in(\mathbf Z/d)^m:
       v_i\equiv0\pmod q,\ \sum_i v_i\equiv0\pmod{e_0}\}.
\]
Its action has $q$ sheet-label orbits.  Connectedness of the composite
cover forces $q=1$, so the full monodromy is $G(d,e_0,m)$.
This proves nonemptiness with the required group, rather than merely
the numerical balance equations.

For a particularly simple explicit family, use $g=0$ and
$r_1=r_{d-1}=n+1$, all other $r_s=0$.  There are $2m-2$ transposition
reflections and $2n+2$ diagonal reflections.  Every Nielsen component
with these data has rank $2n-1$; at $n=5$ it has rank $9$, whereas the
sharp maximum for $G(d,1,m)$ is $14$.
Here the cyclic codifferential can also be checked without Green's
lemma: it is
\[
 H^0(\mathbf P^1,\mathcal O(n-1))^{\otimes2}
 \longrightarrow H^0(\mathbf P^1,\mathcal O(2n-2)),
\]
and every monomial of degree at most $2n-2$ is a product of two
monomials of degree at most $n-1$.

\paragraph{Exact scope left after HD-I12.}
For reflection-only branch data and $n=5$, the cases $g=0,1,2,3$
are now completely determined, with ranks $9,10,11,12$, respectively,
on every component.  HD-I06--HD-I10 already give the sharp maxima
$15$ and $14$ for the remaining range.  HD-I12 alone does not determine
every component at $g=n-1,n,n+1$, or every component with nonreflection
local monodromies.  In those cases the exact composite differential of
HD-I02 and HD-I05 remains the stated componentwise rank formula.

\paragraph{Checks.}
GPT-6 Astra Pro, inherited effort, agent \texttt{/root/imprimitive},
read this whole proof on 21 November 2026, checking the colored-cycle
counts, the precise $H^0$-lemma hypotheses, the Cech obstruction map,
the preservation of the Nielsen component, and the full group
generation.  No numerical test is being substituted for these
universal assertions.  An independent whole-proof reading has been
requested.
